\documentclass[10pt,a4paper]{article}

\usepackage[T1]{fontenc}
\usepackage[utf8]{inputenc}
\usepackage[ngerman]{babel}
\usepackage[a4paper,margin=14mm]{geometry}
\usepackage{lmodern}
\usepackage[table]{xcolor}
\usepackage{tabularx}
\usepackage{array}
\usepackage{microtype}
\usepackage[most]{tcolorbox}

\pagestyle{empty}
\setlength{\parindent}{0pt}
\setlength{\parskip}{0pt}
\renewcommand{\familydefault}{\sfdefault}
\linespread{1.12}\selectfont
\renewcommand{\arraystretch}{1.35}

\definecolor{Navy}{HTML}{17324D}
\definecolor{Blue}{HTML}{3B6EA5}
\definecolor{Sky}{HTML}{EAF3FA}
\definecolor{Mint}{HTML}{E8F5EF}
\definecolor{Gold}{HTML}{F3B33D}
\definecolor{Ink}{HTML}{1E2935}
\definecolor{Muted}{HTML}{617080}
\definecolor{Line}{HTML}{D8E1E8}
\definecolor{CodeBg}{HTML}{F5F7F9}

\color{Ink}
\newcommand{\sectiontitle}[1]{%
  \vspace{4mm}{\color{Blue}\bfseries\Large #1}\par
  \vspace{1mm}{\color{Blue}\rule{\linewidth}{0.7pt}}\par\vspace{2mm}}
\newcommand{\sqlkw}[1]{\textcolor{Blue}{\bfseries #1}}
\newcommand{\sqlbox}[1]{%
  \begin{tcolorbox}[colback=CodeBg,colframe=Line,arc=1.3mm,boxrule=.55pt,
    left=3mm,right=3mm,top=2mm,bottom=2mm,boxsep=0pt]
  {\ttfamily\fontsize{10.2}{14}\selectfont #1}
  \end{tcolorbox}}
\newcommand{\resulttable}[2]{%
  \begin{center}\fontsize{9.5}{12}\selectfont
  \begin{tabular}{|>{\ttfamily\bfseries}c|}\hline
  \rowcolor{Sky}#1\\\hline
  #2\\\hline
  \end{tabular}
  \end{center}}
\newcommand{\pageheader}{%
  \colorbox{Navy}{\parbox{\dimexpr\linewidth-2\fboxsep\relax}{%
    \vspace{3mm}\color{white}\begin{tabularx}{\linewidth}{@{}Xr@{}}
    {\small\bfseries INFORMATIK \textperiodcentered{} KLASSE 10} &
    \colorbox{Gold}{\color{Navy}\small\bfseries\strut LÖSUNGEN}
    \end{tabularx}
    \vspace{2.5mm}{\fontsize{18}{22}\selectfont\bfseries
    Aufgabe 5 -- SQL-Grundlagen wiederholen}\par
    \vspace{1mm}{\large Sicherungsblatt}\vspace{3mm}}}}
\newcommand{\pagefooter}[1]{%
  \vfill{\color{Line}\rule{\linewidth}{.5pt}}\par\vspace{1.5mm}
  \begin{tabularx}{\linewidth}{@{}Xc r@{}}
  {\small\color{Muted}Klasse 10 \textperiodcentered{} Datenbanken} &
  {\small\color{Muted}#1 / 2} &
  {\small\color{Muted}Sicherungsblatt \textperiodcentered{} Aufgabe 5}
  \end{tabularx}}
\newcommand{\merkebox}[1]{%
  \begin{tcolorbox}[colback=Mint,colframe=Blue,arc=2mm,boxrule=.8pt,
    left=4mm,right=4mm,top=3mm,bottom=3mm]
  \textcolor{Blue}{\bfseries Merke}\par\vspace{1.5mm}
  \fontsize{10.5}{14}\selectfont #1
  \end{tcolorbox}}

\begin{document}

\pageheader

\vspace{3mm}
\colorbox{Sky}{\parbox{\dimexpr\linewidth-2\fboxsep\relax}{\vspace{2.5mm}
\textcolor{Blue}{\bfseries Ziel:} Einfache Abfragen lesen und Bedingungen passend verbinden.
\vspace{2.5mm}}}

\sectiontitle{1 \textperiodcentered{} Aufbau einer einfachen Abfrage}
\sqlbox{\sqlkw{SELECT} username, birthday\par
\sqlkw{FROM} users\par
\sqlkw{WHERE} city != 'Berlin';}

\vspace{2mm}
\begin{tabularx}{\linewidth}{@{}>{\ttfamily\bfseries\color{Blue}}l X@{}}
SELECT & legt die Ergebnisspalten fest.\\
FROM & nennt die Tabelle.\\
WHERE & wählt passende Datensätze aus.
\end{tabularx}

\sectiontitle{2 \textperiodcentered{} Mehrere Bedingungen verbinden}
\begin{minipage}[t]{0.48\linewidth}
\textcolor{Blue}{\bfseries \texttt{AND}: beide Bedingungen}\par\vspace{1mm}
\sqlbox{\sqlkw{SELECT} *\par
\sqlkw{FROM} users\par
\sqlkw{WHERE} gender = 'female'\par
\quad \sqlkw{AND} city = 'Leipzig'}
\end{minipage}\hfill
\begin{minipage}[t]{0.48\linewidth}
\textcolor{Blue}{\bfseries \texttt{LIKE}: Name beginnt mit B}\par\vspace{1mm}
\sqlbox{\sqlkw{SELECT} *\par
\sqlkw{FROM} users\par
\sqlkw{WHERE} name \sqlkw{LIKE} 'B\%'}
\texttt{\%} steht für beliebig viele Zeichen.
\end{minipage}

\vspace{4mm}
\textcolor{Blue}{\bfseries \texttt{OR}: eine geklammerte Gruppe genügt}\par\vspace{1mm}
\sqlbox{\sqlkw{SELECT} name\par
\sqlkw{FROM} users\par
\sqlkw{WHERE} (gender = 'male' \sqlkw{AND} centimeters > 165)\par
\quad \sqlkw{OR} (gender = 'female' \sqlkw{AND} centimeters > 160)}

\vspace{3mm}
\merkebox{\texttt{SELECT} bestimmt die Ergebnisspalten, \texttt{FROM} die Tabelle
und \texttt{WHERE} die ausgewählten Datensätze.\par\vspace{2mm}
Textwerte stehen in einfachen Anführungszeichen.}

\pagefooter{1}

\newpage
\pageheader

\sectiontitle{3 \textperiodcentered{} Zusammenfassen, benennen und begrenzen}
\begin{tcolorbox}[colback=white,colframe=Line,arc=1.5mm,boxrule=.6pt,
  left=4mm,right=4mm,top=3mm,bottom=3mm]
\begin{tabularx}{\linewidth}{@{}>{\ttfamily\bfseries\color{Blue}}l X
  >{\ttfamily\bfseries\color{Blue}}l X@{}}
COUNT(*) & Anzahl & MIN(attribut) & kleinster Wert\\
MAX(attribut) & größter Wert & AVG(attribut) & Mittelwert\\
SUM(attribut) & Summe & AS & Spalte umbenennen
\end{tabularx}
\end{tcolorbox}

\vspace{4mm}
\begin{minipage}[t]{0.48\linewidth}
\textcolor{Blue}{\bfseries Anzahl aus München}\par\vspace{1mm}
\sqlbox{\sqlkw{SELECT} \sqlkw{COUNT}(*) \sqlkw{AS} anzahl\par
\sqlkw{FROM} users\par
\sqlkw{WHERE} city = 'München'}
\resulttable{anzahl}{24}
\end{minipage}\hfill
\begin{minipage}[t]{0.48\linewidth}
\textcolor{Blue}{\bfseries Kleinste Körpergröße}\par\vspace{1mm}
\sqlbox{\sqlkw{SELECT} \sqlkw{MIN}(centimeters)\par
\quad \sqlkw{AS} kleinste\_groesse\par
\sqlkw{FROM} users;}
\resulttable{kleinste\_groesse}{152}
\end{minipage}

\vspace{6mm}
\textcolor{Blue}{\bfseries Zuletzt registrierter Datensatz}\par\vspace{1mm}
\sqlbox{\sqlkw{SELECT} *\par
\sqlkw{FROM} users\par
\sqlkw{ORDER BY} created\_at \sqlkw{DESC}\par
\sqlkw{LIMIT} 1;}

\vspace{2mm}
\begin{tabularx}{\linewidth}{@{}>{\ttfamily\bfseries\color{Blue}}l X@{}}
DESC & sortiert absteigend.\\
LIMIT 1 & lässt nur die erste Zeile übrig.
\end{tabularx}

\vspace{8mm}
\merkebox{\texttt{AND} verlangt beide Bedingungen; bei \texttt{OR} genügt eine
geklammerte Gruppe.\par\vspace{2mm}
Aggregatfunktionen liefern aus vielen Werten ein Ergebnis; \texttt{AS} benennt
die Ergebnisspalte.\par\vspace{2mm}
\texttt{ORDER BY ... DESC} sortiert absteigend; \texttt{LIMIT 1} zeigt nur die
erste Ergebniszeile.}

\pagefooter{2}
\end{document}
